\section{Sim-to-Real Validity Framework}
\label{sec:framework}

\paragraph{Setup.}
Figure~\ref{fig:protocol} summarises the protocol. A \emph{package} is an A/B test containing two
or more copy \emph{variants} shown to the same real audience; each variant $v$ has an observed
click-through rate $\mathrm{CTR}(v)$ and a simulator-predicted score $\texttt{pred\_score}(v)$
(Section~\ref{sec:system}).
% Sim-to-real protocol diagram (self-contained TikZ; no external data).
% figure* spans both columns; resizebox keeps the horizontal pipeline within \textwidth.
\begin{figure*}[t]
  \centering
  \resizebox{0.92\textwidth}{!}{%
  \begin{tikzpicture}[
    font=\footnotesize,
    box/.style={draw, rounded corners, align=center, inner sep=5pt, minimum height=10mm, minimum width=22mm},
    data/.style={box, fill=green!8},
    proc/.style={box, fill=blue!6},
    res/.style={box, fill=gray!8},
    >=Stealth, node distance=8mm and 14mm,
  ]
    \node[data] (arch) {Upworthy A/B archive\\\scriptsize(headline variants,\\\scriptsize real click-through)};
    \node[proc, right=of arch] (filt) {Reliability filter\\\scriptsize(significant winner,\\\scriptsize two-proportion $z$)};
    \node[data, right=of filt] (rel) {Reliable\\packages};
    \node[proc, above right=4mm and 16mm of rel] (panel) {Persona panel\\\scriptsize(10, Pew-grounded)};
    \node[proc, below right=4mm and 16mm of rel] (base) {No-persona\\\scriptsize zero-shot baseline};
    \coordinate (fork) at ($(panel.east)!0.5!(base.east)$);
    \node[res, right=12mm of fork, anchor=west] (score) {Click-intent\\\scriptsize$\to$ within-package\\\scriptsize ranking};
    \node[res, right=of score] (val) {Validity vs.\ real CTR\\\scriptsize($\tau$, $\rho$, top-1;\\\scriptsize bootstrap CI)};

    \draw[->] (arch) -- (filt);
    \draw[->] (filt) -- (rel);
    \draw[->] (rel.east) -- (panel.west);
    \draw[->] (rel.east) -- (base.west);
    \draw[->] (panel.east) -- (score.north west);
    \draw[->] (base.east) -- (score.south west);
    \draw[->] (score) -- (val);
  \end{tikzpicture}%
  }
  \caption{The sim-to-real validity protocol. Real A/B tests are filtered to those with a
  statistically reliable winner; the persona panel and a no-persona baseline each score the
  variants, which are ranked within each package and compared to the real click-through ordering.}
  \label{fig:protocol}
\end{figure*}
 Because variants
within a package share an audience, the package defines a controlled comparison.

\paragraph{Why ranking, not absolute CTR.}
We evaluate \emph{within-package ranking} rather than absolute CTR for two reasons. First, CTR
levels vary across packages for reasons unrelated to copy (topic, timing, audience), so only the
within-package contrast is attributable to the copy. Second, the simulator's score and the
benchmark outcome are different constructs (intent vs.\ engagement; Section~\ref{sec:discussion}),
so a ranking criterion is robust to scale and calibration differences. Concretely, a useful
copy simulator need not predict the CTR number---it needs to order candidate variants the way
the real audience does.

\paragraph{Metrics.}
For each package we compare the CTR ordering and the \texttt{pred\_score} ordering with: top-1
winner accuracy (did the predicted best variant match the real best), Kendall's
$\tau$~\citep{kendall1938}, and Spearman's $\rho$~\citep{spearman1904}. Package-level values are
aggregated across packages with bootstrap confidence intervals (Section~\ref{sec:setup}).
